\documentclass[11pt,leqno]{amsart}
\usepackage[letterpaper,margin=1in]{geometry}
\usepackage[T1]{fontenc}
\usepackage{lmodern}
\usepackage{amsmath,amssymb,amsthm,mathtools}
\usepackage{microtype}
\linespread{0.985}
\usepackage{xurl}
\usepackage[colorlinks=true,linkcolor=blue,citecolor=blue,urlcolor=blue]{hyperref}
\title[Erd\H{o}s Problem 12]{The three parts of Erd\H{o}s Problem 12:\\ counting functions and reciprocal sums}
\author[W. Howe]{Will Howe}
\author[G. Huffman]{Grant Huffman}
\author[J. Kohlmeyer]{Jensen Kohlmeyer}
\author[L. Kruer]{Liam Kruer}
\thanks{\textit{E-mail.} \href{mailto:howew@purdue.edu}{\texttt{howew@purdue.edu}} (W. Howe), \href{mailto:huffma97@purdue.edu}{\texttt{huffma97@purdue.edu}} (G. Huffman), \href{mailto:jenwin@purdue.edu}{\texttt{jenwin@purdue.edu}} (J. Kohlmeyer), \href{mailto:kruerl@purdue.edu}{\texttt{kruerl@purdue.edu}} (L. Kruer).}
\date{November 5, 2026}
\subjclass[2020]{11B05, 11N35, 11N36}
\keywords{Erd\H{o}s problems, reciprocal sums, divisibility, sieve methods, primitive sets}
\newtheorem{theorem}{Theorem}[section]
\newtheorem{lemma}[theorem]{Lemma}
\newtheorem{proposition}[theorem]{Proposition}
\newtheorem{corollary}[theorem]{Corollary}
\theoremstyle{definition}
\newtheorem{definition}[theorem]{Definition}
\theoremstyle{remark}
\newtheorem{remark}[theorem]{Remark}
\numberwithin{equation}{section}
\newcommand{\N}{\mathbb N}
\newcommand{\Z}{\mathbb Z}
\newcommand{\T}{\mathbb R/\mathbb Z}
\newcommand{\E}{\mathbb E}
\newcommand{\Prob}{\mathbb P}
\newcommand{\ind}{\mathbf 1}
\newcommand{\hrec}{\mathcal H}
\newcommand{\Pmax}{P^{+}}
\newcommand{\norm}[1]{\left\lVert #1\right\rVert}
\newcommand{\abs}[1]{\left\lvert #1\right\rvert}
\newcommand{\e}{\mathrm e}
\newcommand{\eps}{\varepsilon}
\hypersetup{
 pdftitle={The three parts of Erdos Problem 12: counting functions and reciprocal sums},
 pdfauthor={Will Howe, Grant Huffman, Jensen Kohlmeyer, Liam Kruer},
 pdfsubject={Erdos Problem 12, parts i, ii and iii},
 pdfkeywords={reciprocal sums, divisibility, Property P, sieve methods}
}
\begin{document}
\begin{abstract}
We give a unified account of the three numbered questions in Erd\H{o}s
Problem~12 for sets in which no element divides the sum of two distinct
larger elements. The positive answer to the square-root counting question
and the negative answer to a universal power-saving bound are prior
AlphaProof Nexus results of Google DeepMind. We include attributed proofs
following its released Lean constructions. For the reciprocal-sum question,
we retain the proof that one absolute constant bounds the sum for every
such set, finite or infinite. The distinction between parameterwise dense
examples and the uniform reciprocal bound is made explicit. A companion
single-file Lean formalization combines the two adapted prior proofs with
the reciprocal-sum proof and checks all three statements together.
\end{abstract}
\maketitle
\section{The three questions and their attribution}\label{sec:introduction}

A set \(A\subseteq\N=\{1,2,\ldots\}\) has \emph{Property P} if there
are no pairwise distinct elements \(a,b,c\in A\) such that
\[
 a<b,\qquad a<c,\qquad a\mid b+c.
\]
In particular, equality of the two larger summands is permitted by the
formulation and does not constitute a forbidden triple. Write
\(A(N)=|A\cap\{1,\ldots,N\}|\).
Erd\H{o}s Problem~12 asks three separate questions about the counting
function and reciprocal sum of such an infinite set.

\begin{theorem}[Part (i), AlphaProof Nexus]\label{thm:part-i}
There is an infinite set \(A\) with Property P such that
\[
 \liminf_{N\to\infty}\frac{A(N)}{\sqrt N}>0.
\]
\end{theorem}

\begin{theorem}[Part (ii), AlphaProof Nexus]\label{thm:part-ii}
There is no absolute \(c>0\) such that every infinite set \(A\) with
Property P satisfies \(A(N)<N^{1-c}\) for infinitely many \(N\).
More precisely, for each \(0<c<1\) there are an infinite set \(A_c\)
with Property P and an integer \(N_c\) such that
\begin{equation}\label{eq:parameterwise-dense}
 A_c(N)\ge N^{1-c}\qquad(N\ge N_c).
\end{equation}
The set and the threshold may depend on \(c\).
\end{theorem}

\begin{theorem}[Part (iii), with a uniform bound]\label{thm:main}
There is an absolute constant \(C\) such that every set
\(A\subseteq\N\) with Property P satisfies
\begin{equation}\label{eq:main}
 \hrec(A):=\sum_{a\in A}\frac1a\le C.
\end{equation}
The same constant applies to finite and infinite sets.
\end{theorem}

The answers to the three numbered questions are therefore, respectively,
\emph{yes}, \emph{no}, and \emph{yes}. The first two answers are prior
results of the Google DeepMind AlphaProof Nexus work of Tsoukalas,
Kovsharov, Shirobokov et al.\ \cite{TKS}. They are included here with
attribution, exposition, and adaptations of the released Lean proofs
\cite{NexusI,NexusII}; they are not contributions first obtained in this
manuscript. The reciprocal-sum argument of Theorem~\ref{thm:main} is the
argument developed in the original manuscript by the four named authors.
The present version brings the three statements into one account and one
checked Lean source file.

The quantifier order in \eqref{eq:parameterwise-dense} matters:
\(\forall c\,\exists A_c\,\exists N_c\). We use precisely this
parameterwise assertion. The construction below does not assert that one
fixed set works for every \(c\). The lower bounds in part (ii) are
compatible with the uniform reciprocal bound in part (iii): the onset
\(N_c\) and the set itself vary with the exponent.

Erd\H{o}s and S\'ark\"ozy \cite{ES} proved that every infinite set with
Property P has density zero, and formulated the reciprocal-sum problem
and its stronger uniform version. For the counting questions,
Section~\ref{sec:counting} follows the congruence-block constructions in
the released AlphaProof Nexus Lean files. Those are the concrete proof
instances incorporated in the companion formalization. The constructions
are presented here directly, rather than by relying on admitted statements
in the Formal Conjectures catalog.

The proof of part (iii) has three main steps. A quantitative
incompatibility between separated blocks of rough elements is combined
with sieve estimates to control prime fibres carrying inherited congruence
restrictions. We then represent reciprocal sums by random prime-factor
prefix weights. An adapted rule for selecting the restrictions gives a
uniform total expected weight. The restriction estimate is a special case
of Bera--Viswanadham \cite{BV}; Section~\ref{sec:restriction} includes a
direct proof in the form needed here. The Mertens prefix weights are
classical; see Alexeev et al.\ \cite{AB}. Earlier work of Plasma AI
\cite{Plasma} contains the reduction to odd primitive sets, restricted
reciprocal-sum estimates, and the problem of controlling accumulated
congruence restrictions. These attributions apply to the reciprocal-sum
argument, independently of the provenance of parts (i) and (ii).

\subsection*{Notation and conventions}
All logarithms are natural unless a base is indicated.
We write \(\e(t)=\exp(2\pi i t)\), and \(\norm{t}_{\T}\) for distance
from a real number to the nearest integer. The constants implicit in
\(\ll\) and \(\asymp\) are absolute unless a dependence is indicated.
An integer is \emph{\(D\)-rough} if no prime factor is at most \(D\),
and \emph{\(D\)-smooth} if every prime factor is at most \(D\).
The integer 1 satisfies both conventions. We write \(\Pmax(g)\) for
the largest prime factor of \(g>1\), with \(\Pmax(1)=1\).
A set is \emph{primitive} if no two distinct members divide one another.
Property P is inherited by subsets and by division by a common divisor.
The paper uses positive integers; the Lean statements use natural numbers,
with \(1/0=0\). The constructions below consist of positive integers.

\section{The counting constructions: parts (i) and (ii)}\label{sec:counting}

This section is an attributed exposition of the released proofs
\cite{NexusI,NexusII}. We isolate the common congruence argument and then
make the two choices of scale. Omitting finitely many initial blocks in
the exposition simplifies the estimates and does not affect the asserted
asymptotic statements. The Lean implementations verify explicit initial
bounds as well.

\begin{lemma}[Congruence separation]\label{lem:congruence-separation}
Let \(B_j\) be finite sets of positive integers, with every element of
\(B_i\) smaller than every element of \(B_j\) when \(i<j\). Suppose that
there are pairwise distinct primes \(q_j>2\) such that elements of
\(B_j\) are \(0\pmod{q_j}\) and \(1\pmod{q_i}\) for every \(i<j\).
If each \(B_j\) lies in an interval \([L_j,U_j]\) with
\(2U_j<3L_j\), then \(\bigcup_j B_j\) has Property P.
\end{lemma}
\begin{proof}
Suppose \(a\in B_i\), \(a<b,c\), and \(a\mid b+c\).
If either larger element lies in a later block, then \(b+c\) is
\(1\) or \(2\pmod{q_i}\), whereas \(a\equiv0\pmod{q_i}\).
This is impossible because \(q_i>2\). Otherwise both larger elements
lie in \(B_i\), and
\(2a<b+c\le2U_i<3L_i\le3a\), which precludes divisibility.
\end{proof}

\subsection{A positive square-root lower limit}

Choose increasing primes \(q_0,q_1,\ldots\), starting with \(q_0=5\),
so that \(q_{j+1}\le2q_j+2\). Bertrand's postulate gives such a choice,
and induction yields \(q_j+2\le7\cdot2^j\). Put
\[
 K_j=729\prod_{i<j}q_i,\qquad P_j=K_jq_j,\qquad M_j=P_j^8,
 \qquad T_j=\left\lfloor\frac{P_{j+1}^4}{2}\right\rfloor.
\]
The Chinese remainder theorem supplies \(r_j\in[0,P_j)\) such that
\(r_j\equiv1\pmod{K_j}\) and \(r_j\equiv0\pmod{q_j}\).
Define
\[
 B_j=\{M_j+r_j+tP_j:0\le t<T_j\}.
\]
Then \(|B_j|=T_j\). Also \(P_{j+1}=P_jq_{j+1}\), while
\(P_j\ge3645\cdot5^j\). Consequently
\[
 \frac{q_{j+1}^4}{P_j^3}\longrightarrow0,\qquad
 \frac{T_jP_j}{M_j}\longrightarrow0.
\]
Thus, after removing finitely many initial blocks, each \(B_j\) lies
in \([M_j,(7/5)M_j]\). The scales are strictly separated because
\(M_{j+1}=M_jq_{j+1}^8\). The residues of \(B_j\) have the form in
Lemma~\ref{lem:congruence-separation}, since \(q_i\mid K_j\) for
\(i<j\). Their union \(A\) therefore has Property P.

For large \(j\),
\[
 T_j=\tfrac12\sqrt{M_{j+1}}+O(1).
\]
If \(N\) falls inside the short interval containing \(B_j\), the
preceding complete block contributes
\(T_{j-1}=\tfrac12\sqrt{M_j}+O(1)\) elements. If \(N\) falls in the
following gap, block \(B_j\) contributes
\(T_j=\tfrac12\sqrt{M_{j+1}}+O(1)\) elements. These two cases give
\[
 A(N)\ge\tfrac13\sqrt N
\]
for every sufficiently large \(N\).
To verify that the lower limit is finite as well, note that \(r_j>0\),
so \(A(M_j)=\sum_{i<j}|B_i|\), up to the omitted finite prefix. Since
\(P_{i+1}\le P_j5^{-(j-i-1)}\),
\[
 A(M_j)\le\frac12 P_j^4\sum_{k\ge0}5^{-4k}<P_j^4=\sqrt{M_j}.
\]
This proves Theorem~\ref{thm:part-i}. In particular, the lower limit is
bounded away from zero and bounded above; this agrees with the real-valued
lower-limit convention used by the Lean statement.

\subsection{Failure of a universal power saving}

Fix \(0<c<1\). Choose distinct odd primes \(p_j\) with
\(p_j\le2^{j+2}\), and put \(Q_j=\prod_{i\le j}p_i\). Then
\(Q_j\le2^{(j+2)^2}\). Choose an integer \(a\ge4\) with \(ac\ge3\),
and define
\[
 M_j=2^{a(j+2a)^2}.
\]
Choose a residue \(s_j\pmod{Q_j}\) that is \(0\pmod{p_j}\) and
\(1\pmod{p_i}\) for \(i<j\), and set
\[
 B_j=\{n\in[10M_j,14M_j]:n\equiv s_j\pmod{Q_j}\}.
\]
The intervals are disjoint and ordered because
\(14M_j<10M_{j+1}\). They satisfy the strict width inequality in
Lemma~\ref{lem:congruence-separation}. Once again their union \(A_c\)
has Property P, and it is infinite because the blocks are nonempty for
all sufficiently large \(j\).

Counting a residue class in an interval gives
\[
 |B_j|\ge\frac{4M_j}{Q_j}-1.
\]
The logarithm to base 2 of \(4M_j/Q_j\) is at least
\(2+a(j+2a)^2-(j+2)^2\), whereas
\[
 \log_2\bigl((14M_{j+1})^{1-c}\bigr)
 =(1-c)\bigl(\log_2 14+a(j+1+2a)^2\bigr).
\]
The difference has leading term \((ac-1)j^2\), whose coefficient is
positive. Therefore, for all sufficiently large \(j\),
\[
 |B_j|\ge(14M_{j+1})^{1-c}.
\]
For \(14M_j\le N<14M_{j+1}\), the complete block \(B_j\) is already
counted in \(A_c(N)\). This proves \eqref{eq:parameterwise-dense}.
If a universal saving \(d>0\) existed, apply the construction with
\(c=\min(d,1/2)\). For sufficiently large \(N\ge1\),
\(A_c(N)\ge N^{1-c}\ge N^{1-d}\), contradicting the required
infinitely many strict inequalities. This proves
Theorem~\ref{thm:part-ii}.


\section{A restriction estimate for sifted integers}\label{sec:restriction}

Let \(\mathcal P\) be any set of primes. Define a multiplicative
function supported on squarefree products of primes in \(\mathcal P\)
by
\[
 h_{\mathcal P}(1)=1,\qquad
 h_{\mathcal P}(d)=\prod_{p\mid d}\frac1{p-1},
\]
and put \(h_{\mathcal P}(d)=0\) for all other positive integers \(d\).
For \(z\ge2\), set
\[
 G_{\mathcal P}(z)=\sum_{d<z}h_{\mathcal P}(d),\qquad
 \mathcal S(N,\mathcal P,z)=
 \{1\le n\le N:p\nmid n\text{ for every }p\in\mathcal P,\ p<z\}.
\]
Using strict inequalities at the sieve cutoff removes endpoint
ambiguities.

\begin{lemma}\label{lem:restriction}
Let \(2\le z\le N^{1/4}\), and let \(\mathcal B\subseteq\T\) be
a finite \(\delta\)-separated set. For every fixed \(q>2\) and every
complex sequence \((a_n)\) supported on
\(\mathcal S(N,\mathcal P,z)\),
\begin{equation}\label{eq:restriction}
 \left(\sum_{\xi\in\mathcal B}
   \left|\sum_{n\le N}a_n\e(n\xi)\right|^q\right)^{1/q}
 \ll_q
 \left(\frac{N+\delta^{-1}}{G_{\mathcal P}(z)}
             \sum_{n\le N}|a_n|^2\right)^{1/2}.
\end{equation}
The constant is uniform in \(\mathcal P\).
\end{lemma}

\begin{proof}
We give the enveloping-sieve argument in the form needed here.
For a parameter \(z_0\ge2\), let
\(\mathcal P'=\mathcal P\cap[z_0,z)\), write \(h'=h_{\mathcal P'}\),
and put \(G'=\sum_{d<z}h'(d)\). For \(p\in\mathcal P'\), define
\[
 u_p(n)=\frac{1-p\ind_{p\mid n}}{p-1},\qquad
 u_d(n)=\prod_{p\mid d}u_p(n).
\]
The mean of \(u_p\) modulo \(p\) is zero. Its nonzero additive
Fourier coefficients all equal \(-1/(p-1)\), and
\[
 u_p(n)^2=h'(p)+(h'(p)-1)u_p(n).
\]
Consequently the nonnegative periodic function
\begin{equation}\label{eq:majorant}
 \beta(n)=\frac1{(G')^2}
       \left(\sum_{\substack{d<z\\h'(d)>0}}u_d(n)\right)^2
\end{equation}
equals 1 on \(\mathcal S(N,\mathcal P,z)\).
Orthogonality and the Chinese remainder theorem show that its
constant Fourier coefficient is \(1/G'\).

A nonconstant Fourier coefficient can occur only at a reduced
fraction \(a/\ell\), where \(\ell<z^2\) is squarefree and supported
on \(\mathcal P'\). Its absolute value is at most
\begin{equation}\label{eq:majorant-coeff}
 \frac{3^{\omega(\ell)}h'(\ell)}{G'}
 \ll \frac1{G'\sqrt{\ell}},
\end{equation}
where \(\omega(\ell)\) is the number of distinct prime factors.
Indeed, in a term \(u_d u_e\), a prime dividing \(\ell\) has at
most three membership patterns in \(d,e\), and contributes a
coefficient of magnitude at most \(h'(p)\). A prime outside
\(\ell\) must occur in both factors to contribute a constant
term, which is \(h'(p)\). Summing those common factors costs
at most \(G'\). Finally, \(3/(p-1)\le p^{-1/2}\) for all
sufficiently large primes, and the finitely many remaining primes
contribute an absolute factor.

Separating the prime factors below \(z_0\) gives
\begin{equation}\label{eq:G-comparison}
 G_{\mathcal P}(z)
 \le G'\prod_{p<z_0}(1-1/p)^{-1}
 \ll G'\log(2z_0).
\end{equation}

For a nonzero vector \(c=(c_\xi)_{\xi\in\mathcal B}\), put
\[
 R=\frac{\norm{c}_1^2}{\norm{c}_2^2}\ge1,\qquad z_0=(2R)^2,
 \qquad \eta=\min\{\delta/4,1/(4z^4)\}.
\]
Choose a nonnegative integrable function \(\psi\) majorizing 1
on \([1,N]\), with Fourier transform supported on
\([-\eta,\eta]\), integral \(O(N+\eta^{-1})\), and decay
\(O((1+t^2)^{-1})\). One construction is an absolute multiple
of the convolution of
\(\ind_{[-1/\eta,N+1+1/\eta]}\) with
\[
 \eta\left(\frac{\sin(\pi\eta t)}{\pi\eta t}\right)^2.
\]
Poisson summation implies
\[
 \sum_{n\in\Z}\psi(n)\e(n\theta)=0
 \quad\text{if }\norm{\theta}_{\T}>\eta,
 \qquad
 \left|\sum_{n\in\Z}\psi(n)\e(n\theta)\right|
 \ll N+\eta^{-1}.
\]

Use \(\psi\beta\) to majorize the dual square sum. The constant
coefficient of \(\beta\) contributes only when \(\xi=\xi'\).
For a fixed ordered pair \((\xi,\xi')\), at most one reduced
nonzero fraction \(a/\ell\), \(\ell<z^2\), can satisfy
\(\norm{\xi-\xi'+a/\ell}_{\T}\le\eta\), since distinct such
fractions have distance at least \(z^{-4}\).
Every nonzero denominator in \eqref{eq:majorant-coeff} is at
least \(z_0\). We obtain
\[
 \sum_{n\in\mathcal S(N,\mathcal P,z)}
      \left|\sum_{\xi\in\mathcal B}c_\xi\e(n\xi)\right|^2
 \ll\frac{N+\eta^{-1}}{G'}
       \left(\norm{c}_2^2+z_0^{-1/2}\norm{c}_1^2\right).
\]
If \(z_0\ge z\), then \(\beta=1\) and the second term is absent.
Since \(z^4\le N\), \(z_0^{1/2}=2R\), and
\eqref{eq:G-comparison} holds, this gives
\begin{equation}\label{eq:dual-restriction}
 \sum_{n\in\mathcal S(N,\mathcal P,z)}
      \left|\sum_{\xi\in\mathcal B}c_\xi\e(n\xi)\right|^2
 \ll \frac{N+\delta^{-1}}{G_{\mathcal P}(z)}
                \log(2R)\norm{c}_2^2.
\end{equation}

Write
\[
 U^2=\frac{N+\delta^{-1}}{G_{\mathcal P}(z)}
                 \sum_n|a_n|^2,\qquad
 F(\xi)=\sum_n a_n\e(n\xi).
\]
The one-frequency case of \eqref{eq:dual-restriction} gives
\(\max_{\xi\in\mathcal B}|F(\xi)|\ll U\).
For \(\mathcal E_\lambda=\{\xi\in\mathcal B:|F(\xi)|\ge\lambda U\}\),
choose unit complex phases realizing the absolute values in the
pairing with \(F\). Cauchy--Schwarz and
\eqref{eq:dual-restriction} imply
\[
 \lambda^2|\mathcal E_\lambda|
      \ll\log(2|\mathcal E_\lambda|)
 \quad\text{whenever }\mathcal E_\lambda\ne\varnothing.
\]
Thus
\(
 |\mathcal E_\lambda|
 \ll\lambda^{-2}\log(2+\lambda^{-1})
\)
on the bounded range where the level set is nonempty.
Integrating this estimate against \(q\lambda^{q-1}\) proves
\eqref{eq:restriction}.
\end{proof}

\begin{lemma}\label{lem:sieve-density}
There are absolute constants \(c_0>0\) and \(D_0\) with the
following properties. If \(D\ge D_0\), \(N\ge D\), and
\(z=\tfrac12\min\{D,N^{1/4}\}\), then
\begin{equation}\label{eq:full-density}
 G_{\{\text{all primes}\}}(z)\gg\log D.
\end{equation}
If \(16\le K\le D^{c_0}\), then
\begin{equation}\label{eq:band-density}
 G_{\{p>4K\}}(z)\gg\frac{\log D}{\log(2K)}.
\end{equation}
\end{lemma}

\begin{proof}
Put \(z_1=z^c\), where \(c>0\) is a sufficiently small absolute
constant. In the Euler product
\(\prod_{4K<p\le z_1}(1+1/(p-1))\), normalize the weight of each
squarefree product to a probability distribution. A prime \(p\)
is then selected with probability \(1/p\). The expected logarithm
of the product is
\[
 \sum_{4K<p\le z_1}\frac{\log p}{p}\ll c\log z.
\]
Markov's inequality shows that at least half the total weight is
carried by products smaller than \(z\), after \(c\) is fixed
sufficiently small. Mertens' estimates give
\[
 \prod_{4K<p\le z_1}(1-1/p)^{-1}
       \asymp\frac{\log z_1}{\log(2K)}.
\]
Choose \(c_0\) small enough that \(z_1>8K\) for all sufficiently
large \(D\). Since \(\log z_1\gg\log D\), this proves
\eqref{eq:band-density}. The same argument with all primes at
most \(z_1\) proves \eqref{eq:full-density}.
\end{proof}

\section{Interaction between separated rough blocks}\label{sec:blocks}

\begin{lemma}\label{lem:blocks}
Let \(D\ge D_0\), let \(X<Y\) be powers of two with \(X\ge D\),
and put \(K=Y/X\). Assume \(16\le K\le D^{c_0}\).
Suppose that \(E\subseteq[X,2X)\) and \(T\subseteq[Y,2Y)\)
are \(D\)-rough subsets of one set with Property P.
With
\[
 U=\log D,\qquad u=\frac{U|E|}{X},\qquad v=\frac{U|T|}{Y},
\]
one has
\begin{equation}\label{eq:blocks}
 \sqrt u\,v\ll K^{-1/8}.
\end{equation}
\end{lemma}

\begin{proof}
The one-frequency case of Lemma~\ref{lem:restriction}, together
with \eqref{eq:full-density}, gives \(u,v\ll1\). We may
assume both sets are nonempty.

Work on \(\Z/M\Z\), where \(M=32Y\), with normalized Fourier
transform
\[
 \widehat w(\xi)=\frac1M\sum_{n\bmod M}w(n)\e(-n\xi/M).
\]
Function-space norms are normalized by \(M^{-1}\), while
Fourier-space \(\ell^q\)-norms use counting measure.
Set
\[
 f=U\ind_E,\qquad t=U\ind_T,\qquad
 \alpha=\frac1M\sum_n f(n)=\frac{u}{32K}.
\]
Let \(k_0=K/4\), \(k_1=4K\), and \(J=[k_0,k_1]\cap\Z\).
Define
\[
 b(n)=\frac1{|J|}
       \sum_{\substack{a\in E,\ k\in J\\2ka=n}}f(a).
\]
All the products in this definition are smaller than \(M\).
Its Fourier transform and mean satisfy
\[
 B(\xi):=\widehat b(\xi)
     =\frac1{|J|}\sum_{k\in J}\widehat f(2k\xi),
 \qquad \frac1M\sum_n b(n)=\alpha.
\]
Every prime factor of \(2k\) is at most \(4K<D\). The
\(D\)-rough part of \(n=2ka\) is therefore exactly \(a\).
In particular the representation determines \(a\) uniquely, and
\begin{equation}\label{eq:b-pointwise}
 b(n)\ll U/K.
\end{equation}
The support of \(b\) avoids all primes in \((4K,D]\).

For \(a\in E\), put
\[
 I_a=[k_0a,k_1a]\cap\Z,\qquad
 b_0(n)=\sum_{a\in E}\frac{f(a)}{|I_a|}
              \ind_{\{n=2\ell:\ell\in I_a\}}.
\]
Then \(b_0\) has mean \(\alpha\), is \(O(\alpha)\) everywhere,
and is at least \(c\alpha\) on the even integers in
\([2Y,4Y]\). Indeed, \(|I_a|\asymp Y\), and this latter
interval is contained in the doubled interval \(2I_a\) for
every \(a\in E\). Write \(B_0=\widehat b_0\).

\subsection*{A uniform Fourier comparison}
We claim that
\begin{equation}\label{eq:fourier-comparison}
 \norm{B-B_0}_{\ell^\infty}
       \ll\alpha u^{-1/2}K^{-2/9}.
\end{equation}
Fix a frequency \(\xi\), and put \(\lambda=2\xi/M\pmod1\).
Throughout this comparison take \(q=9/4\).

First suppose that
\(\norm{h\lambda}_{\T}\ge1/(100X)\) for every
\(1\le h\le4K\). The frequencies \(k\lambda\), \(k\in J\),
are \(1/(100X)\)-separated. Lemma~\ref{lem:restriction},
applied on \([1,2X]\) and normalized by \(U/M\), gives an
\(\ell^q\)-norm \(O(\sqrt u/K)\) on these frequencies.
H\"older's inequality then yields
\[
 |B(\xi)|\ll\frac{\sqrt u}{K}K^{-1/q}.
\]
The geometric average over \(I_a\) also gives
\(|B_0(\xi)|\ll\alpha/K\).

Otherwise, for some \(h\le4K\), reduce the nearby rational
approximation to obtain
\[
 \lambda=\frac rd+\beta\pmod1,\qquad
 (r,d)=1,\qquad d\le4K,\qquad
 |\beta|\le\frac1{100dX}.
\]
If \(2\le d\le K^{1/2}\), every \(a\in E\) is a unit modulo
\(d\), so
\[
 \norm{a\lambda}_{\T}
     \ge\frac1d-2X|\beta|\ge\frac1{2d}.
\]
The geometric average over \(J\) therefore gives
\(|B(\xi)|\ll\alpha d/K\ll\alpha K^{-1/2}\).

If \(d>K^{1/2}\), divide \(J\) into consecutive chunks of
length at most \(d\). Within a chunk the frequencies \(k\lambda\)
are separated by at least \(1/(2d)\), since \(d\le4K\ll X\).
Apply Lemma~\ref{lem:restriction} to each chunk and then
average over \(J\). This gives
\[
 |B(\xi)|\ll\frac{\sqrt u}{K}d^{-1/q}
            \ll\frac{\sqrt u}{K}K^{-1/(2q)}.
\]
In both cases with \(d\ge2\), the interval average defining
\(B_0\) is bounded by
\[
 |B_0(\xi)|\ll\frac{\alpha d}{Y}
                 \ll\frac{\alpha K}{Y}.
\]
This term is smaller than the right side of
\eqref{eq:fourier-comparison}, after decreasing \(c_0\) if needed.

It remains to consider \(d=1\). Use the lift of \(\lambda\)
with \(|\lambda|\le1/(100X)\), and put
\(\rho=M|\lambda|\). If \(\rho\ge K^{1/4}\), geometric
summation gives
\[
 |B(\xi)|+|B_0(\xi)|\ll\alpha/\rho
                           \ll\alpha K^{-1/4}.
\]
If \(\rho<K^{1/4}\), couple the uniform measure on \(I_a\)
to the grid \(\{ka:k\in J\}\) by intervals
\([ka,(k+1)a)\). The endpoint adjustment has total variation
\(O(1/K)\), and the phase changes by
\(O(a|\lambda|)=O(\rho/K)\) inside each interval. Hence
\[
 |B(\xi)-B_0(\xi)|\ll\alpha(1+\rho)/K.
\]
These bounds prove \eqref{eq:fourier-comparison}, since
\(1/(2q)=2/9\) and \(u\ll1\).

\subsection*{The forbidden correlation}
By \eqref{eq:b-pointwise},
\[
 \norm b_{L^2(\Z/M\Z)}
       \ll\alpha\sqrt{U/u}.
\]
Apply Lemma~\ref{lem:restriction} to \(b\), using the
prime band \((4K,D]\), the full \(1/M\)-spaced frequency
grid, and \eqref{eq:band-density}. The result is
\[
 \norm B_{\ell^q}\ll\alpha u^{-1/2}\sqrt{\log(2K)}.
\]
Since \(b_0\ll\alpha\), Parseval gives
\(\norm{B_0}_{\ell^2}\ll\alpha\), and thus
\(\norm{B_0}_{\ell^q}\ll\alpha\).
Interpolating with \eqref{eq:fourier-comparison}, with
\(q/9=1/4\), gives
\begin{equation}\label{eq:fourier-nine}
 \norm{B-B_0}_{\ell^9}
 \ll\alpha u^{-1/2}K^{-1/6}\log(2K)^{1/8}
 \ll\alpha u^{-1/2}K^{-1/8}.
\end{equation}
The full-prime restriction estimate likewise gives
\(\norm{\widehat t}_{\ell^q}\ll\sqrt v\).
Fourier inversion and H\"older's inequality, with exponents
\(9\) and \(9/8\), now show that replacing \(b\) by \(b_0\) in
\[
 \E_{x,y\bmod M}t(x)t(y)b(x+y)
\]
changes the expression by at most
\(C\alpha v u^{-1/2}K^{-1/8}\).

Every element of \(T\) is odd. Thus \(x+y\) is an even
integer in \([2Y,4Y]\) whenever \(x,y\in T\), and the expression
with \(b_0\) is at least \(c\alpha v^2\).
The expression with \(b\) vanishes. A nonzero term would give
the integer equality \(x+y=2ka\) with \(x,y>a\) and \(a,x,y\)
in the underlying set. Distinct \(x,y\) contradict Property P.
If \(x=y\), then \(x=ka\), which is impossible because \(k\ge2\)
is \(D\)-smooth and \(x\) is \(D\)-rough. Therefore
\(
 c\alpha v^2\le C\alpha v u^{-1/2}K^{-1/8}
\),
which proves \eqref{eq:blocks}.
\end{proof}

\section{Prime sets with inherited congruence restrictions}\label{sec:prime-tail}

We use the analytic large sieve and the Brun--Titchmarsh
inequality in their classical forms \cite{MV,MVbook}. In
particular, if a set of primes occupies at most
\(\vartheta\varphi(m)\) reduced residue classes modulo \(m\),
then, for a fixed \(a>1\), \(x\ge m^a\), \(x\ge2\), and
\(y\ge x\), partial summation gives
\begin{equation}\label{eq:BT-mass}
 \sum_{\substack{p\in Q\\x\le p<y}}\frac1p
 \ll_a \vartheta
       \left(1+\log\frac{\log y}{\log x}\right).
\end{equation}
Indeed, Brun--Titchmarsh bounds the counting function by
\(O_a(\vartheta t/\log t)\) for \(t\ge x\).

\begin{lemma}\label{lem:prime-uniform}
Every set \(Q\) of primes with Property P satisfies
\(\hrec(Q)\ll1\), with one absolute constant.
\end{lemma}

\begin{proof}
Remove the possible prime 2, which costs at most \(1/2\).
If the remaining set is nonempty, let \(q\) be its least
element. The reciprocal sum over \([q,q^2]\) is \(O(1)\)
by the ordinary upper bound for primes and partial summation.

Enumerate the primes of \(Q\) greater than \(q^2\) as
\(p_1<p_2<\cdots\), and set
\[
 L_0=q,\qquad L_i=qp_1\cdots p_i,\qquad
 x_i=\max\{L_i^2,p_{i+1}\},\qquad y_i=L_{i+1}^2.
\]
For every earlier prime modulus \(r\in\{q,p_1,\ldots,p_i\}\),
the later primes occupy at most one residue in each pair
\(\{a,-a\}\) of nonzero residues modulo \(r\). Otherwise
two distinct later primes have sum divisible by \(r\).
The Chinese remainder theorem therefore gives a fraction
at most \(2^{-i-1}\) of the reduced classes modulo \(L_i\).
Moreover \(\log y_i\le3\log x_i\).
Equation~\eqref{eq:BT-mass} bounds the mass in
\([x_i,y_i)\) by \(O(2^{-i})\).
These intervals cover all remaining primes of \(Q\).
For a finite list, its last prime is smaller than the final
product squared, so the same argument applies.
The resulting geometric series proves the claim.
\end{proof}

\begin{lemma}\label{lem:prime-tail}
There are absolute constants \(A_0\ge3\) and \(C\) with the
following property. Let \(Q\subseteq[Y,\infty)\) be a set of
primes with Property P. Let \(q_1,\ldots,q_k<Y\) be distinct
odd primes, and set
\[
 M=q_1\cdots q_k,\qquad\delta=2^{-k},
\]
with \(M=1\) if \(k=0\). Suppose \(Y\ge M^{A_0}\) and,
for each \(q_i\), the residues of \(Q\) occupy at most one
member of every pair of opposite nonzero classes modulo \(q_i\).
Then
\begin{equation}\label{eq:prime-tail}
 \hrec(Q)\le C\frac{\delta(1+k)}{\log(2Y)}.
\end{equation}
\end{lemma}

\begin{proof}
For bounded \(Y\), the condition \(M^{A_0}\le Y\) bounds
\(M\) and \(k\). Lemma~\ref{lem:prime-uniform} then suffices.
We henceforth take \(Y\) larger than a fixed absolute threshold.
The choice of \(A_0\), made below in terms of \(c_0\), is fixed
once and for all.

The Chinese remainder theorem puts all primes of \(Q\) in at
most \(\delta\varphi(M)\) reduced residue classes modulo \(M\).
For any dyadic \(V\ge Y/2\), Brun--Titchmarsh gives
\begin{equation}\label{eq:bin-density}
 |Q\cap[V,2V)|\ll\frac{\delta V}{\log V}.
\end{equation}
Here \(M\le Y^{1/A_0}\), and \(A_0\ge3\).
For \(M=1\), use the ordinary prime upper bound.
We will also use
\begin{equation}\label{eq:delta-lower}
 \delta\ge M^{-1}\ge Y^{-1/A_0}.
\end{equation}

Call the dyadic bin at \(V=2^j\) \emph{light} if its
cardinality is at most \(\delta V/j^2\). All light bins
together contribute \(O(\delta/\log Y)\) to \(\hrec(Q)\).
If there is no other bin, the proof is complete.

Let \([U,2U)\), \(U=2^J\), be the first bin that is not
light, and put \(n=|Q\cap[U,2U)|>\delta U/J^2\).
Choose
\[
 s=\lceil8\log_2J\rceil
\]
distinct primes from that bin. This is possible for all
sufficiently large \(Y\), by \eqref{eq:delta-lower}.
Let \(L\) be \(M\) times their product, and set
\[
 V_0=U J^8\delta^{-8}.
\]
Choose \(A_0\) sufficiently large that, uniformly for
\(U\ge Y/2\),
\begin{equation}\label{eq:near-range}
 4V_0/U\le(U/2)^{c_0}
\end{equation}
for all sufficiently large \(Y\). This follows from
\eqref{eq:delta-lower}. The same range gives \(L>V_0\).

\subsection*{The near range}
Take \(D=U/2\). In dyadic bins meeting \([U,V_0]\), every
prime is \(D\)-rough. The normalized count
\[
 w_V=\frac{\log D}{V}|Q\cap[V,2V)|
\]
is \(O(\delta)\), by \eqref{eq:bin-density}.
If two such counts are at least \(w>0\), then
Lemma~\ref{lem:blocks} bounds their ratio of scales by
\(O(w^{-12})\), unless the scales differ by a factor less
than 16. The hypotheses of that lemma hold throughout
the range by \eqref{eq:near-range}.
It follows that
\[
 \#\{V:w_V\ge w\}\ll1+\log(C/w).
\]
Integrating this bound over \(0<w\ll\delta\) gives
\begin{equation}\label{eq:near-mass}
 \hrec(Q\cap[U,V_0])
 \ll\frac{\delta(1+\log(1/\delta))}{\log U}.
\end{equation}

\subsection*{The middle range}
Let \(T=Q\cap[V,2V)\), where \(V>2U\).
For every prime \(p\in Q\cap[U,2U)\), the set \(T\)
occupies at most \((p-1)/2\) classes modulo \(p\).
Orthogonality and Cauchy--Schwarz give
\[
 \sum_{j=1}^{p-1}
      \left|\sum_{t\in T}\e(jt/p)\right|^2
 \ge |T|^2.
\]
The reduced fractions \(j/p\) belonging to different such
primes are distinct and have spacing at least \(1/(4U^2)\).
The analytic large sieve therefore gives
\[
 n|T|^2\ll(V+U^2)|T|,
 \qquad
 \hrec(T)\ll\frac1n+\frac{U^2}{Vn}.
\]
Summing over bins meeting \([V_0,L]\), including the
at most two boundary bins, yields
\begin{equation}\label{eq:middle-mass}
 \hrec(Q\cap[V_0,L])
 \ll\frac{sJ^3}{\delta U}+\frac{\delta^7}{J^6}
 \ll\frac{\delta}{J}.
\end{equation}
For the first inequality use \(\log L=O(sJ)\) and
\(n>\delta U/J^2\). For the second, divide by
\(\delta/J\) and use
\[
 \frac{sJ^4}{\delta^2U}+\frac{\delta^6}{J^5}=o(1),
\]
uniformly in the permitted range, by \eqref{eq:delta-lower}.

\subsection*{The far range}
On \((L,L^2]\), use the external \(q_i\) and the
\(\lfloor s/2\rfloor\) smallest selected internal primes as
moduli. Their product is at most \(\sqrt{ML}\), hence at
most \(L^{2/3}\) for sufficiently large \(Y\). The occupied
fraction of reduced classes is at most
\(\delta2^{-\lfloor s/2\rfloor}\). Equation~\eqref{eq:BT-mass}
therefore bounds the mass in this range by
\(O(\delta2^{-\lfloor s/2\rfloor})\).

Enumerate the primes of \(Q\) above \(L^2\) as
\(p_1<p_2<\cdots\), and now put \(L_i=Lp_1\cdots p_i\),
with \(L_0=L\). The primes in \([L_i^2,L_{i+1}^2)\) are
at least \(p_{i+1}\), and occupy at most the fraction
\(\delta2^{-s-i}\) of the reduced classes modulo \(L_i\).
With
\[
 x_i=\max\{L_i^2,p_{i+1}\},\qquad y_i=L_{i+1}^2,
\]
one has \(\log y_i\le3\log x_i\). Another application
of \eqref{eq:BT-mass} gives mass \(O(\delta2^{-s-i})\).
As in Lemma~\ref{lem:prime-uniform}, these intervals also
cover every term of a finite list. Altogether the far
range contributes
\[
 O(\delta2^{-s/2})=O(\delta/J).
\]
Combine this with the light-bin bound, \eqref{eq:near-mass},
and \eqref{eq:middle-mass}. Since
\(\log(1/\delta)=k\log2\) and \(\log U\gg\log Y\),
we obtain \eqref{eq:prime-tail}.
\end{proof}

\section{Prime-factor prefixes}\label{sec:prefixes}

Let \(A\) be an odd primitive set with Property P.
If \(1\in A\), then \(A=\{1\}\), so suppose \(1\notin A\).
For each positive integer \(g\), define
\[
 Q_g=\{p\text{ prime}:p\ge\Pmax(g),\ gp\in A\}.
\]
Each \(Q_g\) has Property P, and all its primes are odd.
Every element of \(A\) has a unique representation \(gp\)
in which \(p\) is its largest prime factor. Hence
\begin{equation}\label{eq:fibre-decomposition}
 \hrec(A)=\sum_{g\ge1}\frac{\hrec(Q_g)}g.
\end{equation}
The term at \(g=1\) is \(O(1)\), by
Lemma~\ref{lem:prime-uniform}.

Independently for all primes \(p\), let \(V_p\) have
geometric distribution
\[
 \Prob(V_p=j)=(1-1/p)p^{-j},\qquad j=0,1,2,\ldots.
\]
Form the ordered list consisting of \(V_p\) copies of \(p\),
with primes in increasing order. A \emph{prefix product}
is the product of an initial finite segment of this list.

\begin{lemma}\label{lem:prefix-probability}
For \(g>1\), the probability that \(g\) is a prefix product is
\begin{equation}\label{eq:prefix-probability}
 \mu_g=\frac1g\prod_{p<\Pmax(g)}(1-1/p)
           \asymp\frac1{g\log(2\Pmax(g))}.
\end{equation}
\end{lemma}

\begin{proof}
Put \(r=\Pmax(g)\). For primes \(p<r\), the valuation
\(V_p\) must equal \(v_p(g)\). At \(r\), the condition is
\(V_r\ge v_r(g)\), which has probability \(r^{-v_r(g)}\).
Independence gives the exact formula, and Mertens' product
estimate gives the comparison.
\end{proof}

Assign to a prefix \(g>1\) the score
\[
 W_g=\log(2\Pmax(g))\,\hrec(Q_g).
\]
Tonelli's theorem and \eqref{eq:prefix-probability} show
that it is enough to prove
\begin{equation}\label{eq:expected-score}
 \E\sum_{\substack{g>1\\g\text{ a prefix}}}W_g\ll1.
\end{equation}
Lemma~\ref{lem:prime-tail} with \(k=0\) gives \(W_g\ll1\).
More precisely, if \(Q_g\ne\varnothing\) and
\(p_g=\min Q_g\), then
\begin{equation}\label{eq:score-basic}
 W_g\ll\frac{\log(2\Pmax(g))}{\log(2p_g)}.
\end{equation}
We reserve \(p_g\) for this least terminal prime and use
\(r=\Pmax(g)\) for the prime coordinate of the prefix.

We may omit scores produced by the second or a later
copy of a prime. Their total expectation is bounded by
\[
 C\sum_p\E(V_p-1)_+
     =C\sum_p\frac1{p(p-1)}<\infty.
\]
The omitted copies remain in the list and in all later
prefix products. Scores at primes below any fixed absolute
threshold also have bounded total expectation.
Thus only the first copy at each sufficiently large
prime coordinate needs to be considered.

\begin{lemma}\label{lem:inheritance}
Suppose \(h\) and \(g\) are prefix products on one path,
\(h\) occurs earlier, \(Q_h,Q_g\ne\varnothing\), and
\(\Pmax(g)>p_h\). Then the residues of \(Q_g\) occupy
at most one class from each opposite pair of nonzero
classes modulo \(p_h\). The terminal primes \(p_h\)
from distinct nonempty prefixes on a path are distinct.
\end{lemma}

\begin{proof}
We have \(h\mid g\). If \(p_h\mid g/h\), then
\(hp_h\in A\) divides every \(gp\in A\), \(p\in Q_g\).
These are distinct elements because \(\Pmax(g)>p_h\),
contrary to primitivity. Thus \((g/h,p_h)=1\).
All primes in \(Q_g\) exceed \(p_h\).
For distinct \(p,p'\in Q_g\), the congruence
\(p+p'\equiv0\pmod{p_h}\) would now imply
\[
 hp_h\mid gp+gp',
\]
with both summands larger than \(hp_h\), contradicting
Property P. This proves the first assertion.
If two distinct nested prefixes \(h_1\mid h_2\) had the
same terminal prime \(p\), then the distinct elements
\(h_1p,h_2p\in A\) would contradict primitivity.
\end{proof}

\section{An adapted summation of prefix scores}\label{sec:scheduling}

We process every prime coordinate \(r\) in increasing order,
including coordinates for which \(V_r=0\). At coordinate
\(r\), reveal \(V_r\), retain all its copies in the prefix
product, and score the first copy if \(V_r\ge1\).
Let \(A_0\) be the exponent in Lemma~\ref{lem:prime-tail}.
The following scheduling rule uses selected terminal
primes as congruence moduli, which we call \emph{guards}.

After \(k\) guards have been activated, write \(M\) for
their product and \(\delta=2^{-k}\). A \emph{ready stage}
starts at the first unprocessed prime at least \(R\),
where \(R\ge M^{A_0}\) is determined by the previously
revealed coordinates. Initially \(k=0\) and \(R\) is a
fixed sufficiently large absolute threshold. A stage
starts with no pending guard.

At a scored prefix \(g\) with \(Q_g\ne\varnothing\),
schedule \(p_g\), retaining the smallest pending prime.
If the newly scheduled minimum equals the current
coordinate, activate it immediately after scoring.
A guard already pending when its coordinate is reached
is activated before revealing that coordinate, so that
the coordinate belongs to the delay described below.
The stage ends on activation.

Every finite pending prime is eventually reached, since
all prime coordinates are processed. A ready stage that
never activates therefore has no nonempty scored prefix.
The guard selected in a completed stage arose from a
previously visited prefix. Lemma~\ref{lem:inheritance}
shows that the old guards impose the required opposite-class
restrictions on every nonempty fibre in that ready stage.
They are distinct and smaller than its current coordinate.

At a nonempty scored prefix \(g\) with \(r=\Pmax(g)\),
Lemma~\ref{lem:prime-tail}, applied with \(Y=p_g\ge r\),
gives
\begin{equation}\label{eq:score-guarded}
 W_g\ll \delta(1+k)\frac{\log r}{\log p_g}.
\end{equation}
The harmless factors 2 in \eqref{eq:score-basic} are
absorbed by the fixed lower threshold for \(r\).

\begin{lemma}\label{lem:ready-cost}
Conditional on the history at the start of a ready
stage with \(k\) active guards, the expected total score
in that stage is \(O(2^{-k}(1+k))\).
\end{lemma}

\begin{proof}
Call a change of the pending minimum \emph{large} if its
logarithm decreases by a factor of at least 2. The first
finite pending value is counted as large.
At consecutive large changes, the ratio
\(\log r/\log p_g\) at least doubles, since \(r\)
increases and the pending logarithm halves.
Every such ratio is at most 1 before activation.
Thus the sum of these ratios over all large changes
is at most 2 on every path. Their total score is bounded
by \(C\delta(1+k)\), by \eqref{eq:score-guarded}.

For all other nonempty scored prefixes, let \(P_{\rm old}\)
be the pending prime known before revealing \(V_r\).
Then
\[
 \frac{\log r}{\log p_g}
       \le2\frac{\log r}{\log P_{\rm old}}.
\]
Interpret the right side as zero when no finite prime
is pending. In that case the first nonempty prefix
has already been assigned to the large-change class.

Both readiness at \(r\) and \(P_{\rm old}\) are measurable
before \(V_r\) is revealed. A first-copy score at \(r\)
can occur only if \(V_r\ge1\), an event of conditional
probability \(1/r\).
If \(P_*\) is the prime eventually activated, every ready
coordinate is at most \(P_*\), and every previous finite
pending value is at least \(P_*\). Consequently, pathwise,
\begin{equation}\label{eq:predictable-budget}
 \sum_{\substack{r\text{ ready}\\P_{\rm old}<\infty}}
       \frac1r\frac{\log r}{\log P_{\rm old}}
 \le\frac1{\log P_*}
       \sum_{\substack{r\le P_*\\r\text{ prime}}}
                         \frac{\log r}{r}
 \ll1.
\end{equation}
For a stage without activation the sum is zero.
Taking conditional expectation before each coordinate,
and then applying Tonelli, proves the same
\(O(\delta(1+k))\) bound for all remaining scores.
\end{proof}

After activating \(P_*\), put \(M_{\rm new}=MP_*\).
Process the remaining coordinates below \(M_{\rm new}^{A_0}\)
as a \emph{delay}, then start the next ready stage at the
first unprocessed prime at least that threshold.
The old guards remain valid throughout the delay.

\begin{lemma}\label{lem:delay-cost}
Conditional on the activation history, the expected
total score in the delay is \(O(2^{-k}(1+k))\).
\end{lemma}

\begin{proof}
Equation~\eqref{eq:score-guarded} bounds every scored
prefix during the delay by \(C\delta(1+k)\), using
\(\log r/\log p_g\le1\) and the old guards.
Moreover \(M^{A_0}\le R\le P_*\), so
\[
 M_{\rm new}^{A_0}
       =(MP_*)^{A_0}\le P_*^{A_0+1}.
\]
Conditional on the activation history, all strictly
later prime coordinates retain their independent
geometric laws. Their expected number of first copies
before the next ready stage is at most
\[
 \sum_{\substack{P_*<r\le P_*^{A_0+1}\\r\text{ prime}}}\frac1r
       =O_{A_0}(1)
\]
by Mertens' estimate. If \(P_*\) was pending before its
coordinate, that coordinate is also in the delay and
contributes at most one additional score. If the guard
was scheduled at its own coordinate, that score was
already counted in the ready stage.
This proves the claim.
\end{proof}

\begin{proposition}\label{prop:odd-primitive}
Every odd primitive set \(A\) with Property P satisfies
\(\hrec(A)\ll1\), with one absolute constant.
\end{proposition}

\begin{proof}
Stage-start and activation histories are measurable at
their respective prime coordinates. The conditional
expectation statements in Lemmas~\ref{lem:ready-cost}
and \ref{lem:delay-cost} can therefore be applied on
each event specifying these coordinates. All such
coordinates are countable, and the rule never conditions
on a future valuation.

Each completed ready stage adds one new guard. Summing
the two lemmas over \(k\ge0\), and restoring the omitted
small-prime and repeated-copy scores, gives
\[
 \E\sum_{\substack{g>1\\g\text{ a prefix}}}W_g
 \ll1+\sum_{k\ge0}2^{-k}(1+k)<\infty.
\]
Equation~\eqref{eq:expected-score}, Lemma~\ref{lem:prefix-probability},
and the decomposition \eqref{eq:fibre-decomposition}
complete the proof.
\end{proof}

\section{Completion of the proof}\label{sec:completion}

\begin{proof}[Proof of Theorem~\ref{thm:main}]
It is enough to bound \(\hrec(B)\) uniformly for finite
sets \(B\) with Property P. Let \(C\subseteq B\) consist
of the elements minimal under divisibility.
Every \(b\in B\setminus C\) has a proper divisor in \(C\).
Assign \(b\) to one such divisor \(c\).
No \(c\) receives two distinct elements, since it
would divide their sum and both would exceed it.
Every assigned \(b\) is at least \(2c\). Hence
\begin{equation}\label{eq:primitive-reduction}
 \hrec(B)\le\frac32\hrec(C).
\end{equation}
The set \(C\) is primitive.

For \(v\ge0\), define
\[
 C_v=\{c/2^v:c\in C,\ v_2(c)=v\}.
\]
Each \(C_v\) is an odd primitive set with Property P.
Proposition~\ref{prop:odd-primitive} supplies one constant
\(C_{\rm odd}\) bounding all their reciprocal sums,
including the possible singleton \(\{1\}\).
Therefore
\[
 \hrec(B)\le\frac32\sum_{v\ge0}2^{-v}\hrec(C_v)
                \le3C_{\rm odd}.
\]
This bound is independent of \(B\). Apply it to the
finite truncations of an arbitrary \(A\), and use
monotone convergence to obtain \eqref{eq:main}.
\end{proof}

\section{An explicit finite prefix bound}\label{sec:finite-prefix}

We give a finite formulation of the argument in
Section~\ref{sec:scheduling}, with explicit bounds for the total expected
score. Fix an odd primitive set \(B\subseteq[2,N]\) with Property P.
In this section, the fibres \(Q_g\) and scores \(W_g\) are defined using
\(B\) in place of \(A\). Independently for each prime \(p\le N\), choose
\(X_p\in\{0,1,\ldots,N+1\}\) with
\[
 \Prob(X_p=j)=(1-1/p)p^{-j}\quad(0\le j<N+1),\qquad
 \Prob(X_p=N+1)=p^{-(N+1)}.
\]
For \(0\le a\le N+1\), summing the geometric series gives
\(\Prob(X_p\ge a)=p^{-a}\). Thus, for every \(2\le g\le N\), the event
that \(g\) is a prefix has exactly the probability in
Lemma~\ref{lem:prefix-probability}. Every expectation in this section is
a finite sum.

Apply the scheduling rule of Section~\ref{sec:scheduling}, processing
integer times \(r=1,\ldots,N\) and changing the state only at prime
coordinates. Charge an activation coordinate using its incoming set
of guards. Let \(k_r\) be the number of guards activated before time
\(r\), and put
\[
 w_k=2^{-k}(1+k),\qquad R_k=2^{-k}(2k+4),\qquad
 a_r=R_{k_r}-R_{k_{r+1}}.
\]
Since \(R_0=4\) and \(R_k-R_{k+1}=w_k\), the quantity \(a_r\) is
\(w_{k_r}\) at an activation and zero otherwise. In particular,
\begin{equation}\label{eq:finite-activation}
 \sum_{r=1}^N a_r\le4.
\end{equation}
Choose \(C_{\rm tail}\) uniformly in the prime-tail estimate so that
\(W_g\le C_{\rm tail}\) and, for a ready first-copy prefix with
\(r=\Pmax(g)\),
\[
 W_g\le C_{\rm tail}w_{k_r}
       \frac{\log(2r)}{\log(2p_g)}
 \le 2C_{\rm tail}w_{k_r}\frac{\log r}{\log p_g}.
\]
Increase the fixed initial threshold if necessary.

\subsection{The ready coordinates}

If the pending terminal prime before time \(r\) is \(P\), define
\[
 U_r=w_{k_r}\frac{\log r}{\log P},\qquad
 T_r=\frac{w_{k_r}}{\log P}
           \sum_{\substack{p<r\\p\ \mathrm{prime}}}\frac{\log p}{p}.
\]
Set both quantities to zero when no prime is pending. A pending prime
is at least the current coordinate. Between activations its value can
only decrease.

Write \(W_r^{(1)}\) for the first-copy score charged at a ready time,
and set it to zero at other times. Then
\begin{equation}\label{eq:finite-ready-charge}
 W_r^{(1)}\le 4C_{\rm tail}
 \left(\ind_{\{X_r\ge1\}}U_r+U_{r+1}-U_r+2a_r\right),
\end{equation}
where the indicator is zero at a composite time. To verify this
inequality, first suppose that no activation occurs. If a nonempty
prefix is scored, the new pending prime is at most its terminal prime
\(p_g\), and consequently
\[
 U_{r+1}\ge w_{k_r}\frac{\log r}{\log p_g}.
\]
On this event the first three terms in parentheses in
\eqref{eq:finite-ready-charge} sum to \(U_{r+1}\). If no score occurs,
\(U_{r+1}\ge U_r\). At an activation, \(U_{r+1}=0\),
\(U_r\le w_{k_r}\), and \(a_r=w_{k_r}\), so the activation term
covers both the lost potential and the score. During a delay there is
no pending prime and these quantities are zero.

The state before time \(r\) depends only on earlier coordinates.
Independence therefore gives, for prime \(r\),
\[
 \E\left[\ind_{\{X_r\ge1\}}U_r\right]=\frac1r\E U_r.
\]
Choose \(C_{\log}>0\) such that
\(\sum_{p\le x}(\log p)/p\le C_{\log}\log x\) for \(x\ge2\).
The weighted prime sum and the monotonicity of the pending minimum
imply
\begin{equation}\label{eq:finite-prime-potential}
 \ind_{\{r\text{ prime}\}}\frac{U_r}{r}
 \le T_{r+1}-T_r+C_{\log}a_r.
\end{equation}
Indeed, without activation the increment of the numerator accounts
for \((\log r)/r\), and a decrease of the denominator increases the
potential. At an activation, the numerator through time \(r\) is at
most \(C_{\log}\log P\), so the lost potential together with the
current summand is at most \(C_{\log}w_{k_r}\).

Every terminal prime of a nonempty fibre is at most \(N\). Hence no
prime is pending after time \(N\), and both potentials begin and end
at zero. Summing \eqref{eq:finite-prime-potential} and using
\eqref{eq:finite-activation} gives
\[
 \sum_{r=1}^N\E\left[\ind_{\{X_r\ge1\}}U_r\right]
 \le4C_{\log}.
\]
Summation of \eqref{eq:finite-ready-charge} now bounds the expected
ready contribution by
\begin{equation}\label{eq:finite-ready-total}
 C_{\rm tail}(32+16C_{\log}).
\end{equation}

\subsection{The delayed coordinates}

Suppose that the latest activation was at a prime \(P\), with \(j\)
guards active before that activation and old product \(M\). Put
\(H=(MP)^{A_0}\). The inequality \(M^{A_0}\le P\) gives
\(H\le P^{A_0+1}\). Mertens' estimate supplies an absolute constant
\(C_{\rm delay}>0\), depending only on the fixed \(A_0\), such that
\[
 \sum_{\substack{P\le p<P^{A_0+1}\\p\ \mathrm{prime}}}\frac1p
 \le C_{\rm delay}
\]
for every possible activation prime \(P\). Define the accumulated
mass after this activation by
\[
 B_r=w_j\sum_{\substack{P\le p<\min(r,H)\\p\ \mathrm{prime}}}\frac1p.
\]
Then \(0\le B_r\le C_{\rm delay}w_j\). For \(k_r\ge1\), so that
\(j=k_r-1\), set
\[
 D_r=C_{\rm delay}R_{k_r}+C_{\rm delay}w_{k_r-1}-B_r.
\]
Before the first activation set \(D_r=C_{\rm delay}R_0\).
These quantities are nonnegative. At a delayed prime coordinate the
increment of \(B_r\) is \(w_j/r\), so \(D_r\) decreases by that amount.
At a ready coordinate without activation it is constant.

At a new activation, the previous remaining term
\(C_{\rm delay}w_{k_r-1}-B_r\) is nonnegative. The new value satisfies
\[
 D_{r+1}\le C_{\rm delay}(R_{k_r+1}+w_{k_r})
           =C_{\rm delay}R_{k_r}\le D_r.
\]
Thus the total predictable delayed cost is at most
\(D_1=4C_{\rm delay}\). Each delayed first-copy score is at most
\(C_{\rm tail}w_j\), and its occurrence has conditional probability
\(1/r\). Its total expected contribution is therefore at most
\begin{equation}\label{eq:finite-delay-total}
 4C_{\rm tail}C_{\rm delay}.
\end{equation}

The finitely many initial coordinates contribute an absolute
constant. The repeated-copy contribution is bounded using
\[
 \sum_p\E(X_p-1)_+
 \le\sum_p\frac1{p(p-1)}<\infty.
\]
Combining these estimates with \eqref{eq:finite-ready-total} and
\eqref{eq:finite-delay-total} yields a bound for
\(\E\sum_{g\text{ a prefix}}W_g\) independent of \(B\) and \(N\).
Lemma~\ref{lem:prefix-probability} and
\eqref{eq:fibre-decomposition} then give the uniform reciprocal bound
for finite odd primitive sets. Section~\ref{sec:completion} extends it
to all sets with Property P.


\section{The single-file formalization and provenance}\label{sec:formalization}

The companion file \texttt{Erdos12AllParts.lean} imports only Mathlib.
It contains the local dependency closure of the reciprocal-sum argument
and the two attributed AlphaProof Nexus proof developments. Its public
interface is as follows.

\begin{center}
\begin{tabular}{p{0.22\textwidth}p{0.69\textwidth}}
\hline
Entry point & Mathematical statement\\
\hline
\texttt{part\_i} & A good infinite set has positive lower limit at square-root scale.\\[3pt]
\texttt{part\_ii} & The universal power-saving assertion is false.\\[3pt]
\texttt{part\_iii} & Every good infinite set has a convergent reciprocal sum.\\[3pt]
\texttt{all\_parts} & The conjunction of the three preceding answers.\\[3pt]
\texttt{parameterwise\_}\newline\texttt{dense} & The set and eventual counting threshold may depend on the exponent.\\[3pt]
\texttt{uniform\_}\newline\texttt{reciprocal\_bound} & One constant bounds the reciprocal sum of every set with Property P.\\
\hline
\end{tabular}
\end{center}
All entry points belong to the namespace \texttt{Erdos12AllParts}.
The definition \texttt{IsGood} retains the quantifiers and strict
inequalities of the original Formal Conjectures statement.

The check uses Lean \texttt{4.35.0-rc2} and Mathlib commit
\nolinkurl{065356127b1dc0016f66b7283ce0ce2c4055aa55}.
The prior proof files were taken from Google DeepMind's
\texttt{alphaproof-nexus-results} repository at commit
\nolinkurl{0647711a71183c1ea492ad60860776617ce1ea88}.
They are distributed under the Apache License, Version 2.0, with
copyright 2025 Google LLC. The companion file retains these notices and
the license. Their original environment used Lean 4.27.0;
the adaptation isolates their namespaces, removes redundant options and
imports, updates the image-cardinality API, and expands the
interval-residue counting lemma into a direct proof.

The analytic formalization also retains the Euler--Maclaurin and
Mertens developments from \texttt{PrimeNumberTheoremAnd}, obtained via
\texttt{chen\_theorem}, and the vendored large-sieve development from
\texttt{goldbach-lean} \cite{PNTCode,GoldbachCode}. Their recorded source
revisions and attributions are preserved in the Lean header and the
validation/provenance files; these components are likewise recorded under
Apache License 2.0 in the supplied project.

The final checks compile the combined source, check each public entry
point, and inspect its axiom dependencies. The only axioms reported are
\texttt{propext}, \texttt{Classical.choice}, and \texttt{Quot.sound}.
The file introduces no admitted theorem and does not obtain the first
two results by importing placeholder catalog declarations. A separate
check compares all three endpoints with the pinned original catalog
statements. The validation receipt records the source digest and commands.

\section*{Acknowledgments and contribution statement}

Parts (i) and (ii), including the constructions and the released formal
proofs adapted here, are credited to the AlphaProof Nexus work of
Google DeepMind and Tsoukalas, Kovsharov, Shirobokov et al.\ \cite{TKS}.
The present manuscript supplies a unified exposition and the recorded
adaptation to the pinned Lean environment. The four named authors retain
responsibility for the reciprocal-sum argument, the integration, and the
final mathematical claims.

ChatGPT and Codex, developed by OpenAI, assisted with the mathematical
development of the reciprocal-sum argument, checking calculations,
literature searches, source integration, and manuscript preparation.
The earlier reductions and analytic ingredients used in part (iii)
are credited in the introduction and body of the proof.

\begin{thebibliography}{99}

\bibitem{AB}
B. Alexeev, K. Barreto, Y. Li, J. D. Lichtman, L. Price,
J. I. Shah, Q. Tang and T. Tao,
\emph{Primitive sets and von Mangoldt chains: Erd\H{o}s Problem
\#1196 and beyond},
preprint (2026),
\href{https://arxiv.org/abs/2605.00301v1}{arXiv:2605.00301v1}.

\bibitem{BV}
T. Bera and G. K. Viswanadham,
\emph{Restriction estimates with sifted integers},
preprint (2026),
\href{https://arxiv.org/abs/2512.21640v3}{arXiv:2512.21640v3}.

\bibitem{ES}
P. Erd\H{o}s and A. S\'ark\"ozy,
\emph{On the divisibility properties of sequences of integers},
Proc. London Math. Soc. (3) \textbf{21} (1970), 97--101.
\href{https://doi.org/10.1112/plms/s3-21.1.97}{doi:10.1112/plms/s3-21.1.97}.

\bibitem{NexusI}
Google DeepMind,
\emph{AlphaProof Nexus results: Erd\H{o}s Problem 12(i)},
Lean source (2026), Apache License, Version 2.0.
Repository commit \nolinkurl{0647711a71183c1ea492ad60860776617ce1ea88}.
\url{https://github.com/google-deepmind/alphaproof-nexus-results/blob/0647711a71183c1ea492ad60860776617ce1ea88/APNOutputs/ErdosProblems/erdos_12.parts.i.lean}.

\bibitem{NexusII}
Google DeepMind,
\emph{AlphaProof Nexus results: Erd\H{o}s Problem 12(ii)},
Lean source (2026), Apache License, Version 2.0.
Repository commit \nolinkurl{0647711a71183c1ea492ad60860776617ce1ea88}.
\url{https://github.com/google-deepmind/alphaproof-nexus-results/blob/0647711a71183c1ea492ad60860776617ce1ea88/APNOutputs/ErdosProblems/erdos_12.parts.ii.lean}.

\bibitem{MV}
H. L. Montgomery and R. C. Vaughan,
\emph{\href{https://doi.org/10.1112/S0025579300004708}{The large sieve}},
Mathematika \textbf{20} (1973), 119--134.

\bibitem{MVbook}
H. L. Montgomery and R. C. Vaughan,
\emph{Multiplicative Number Theory I. Classical Theory},
Cambridge Studies in Advanced Mathematics, vol.~97,
Cambridge University Press, Cambridge, 2007.

\bibitem{Plasma}
Plasma AI,
\emph{Composite-modulus large sieve for Property P},
research notes, Erd\H{o}s Problems (2026),
\href{https://erdosproblems.ai/wiki/research/leads/property_p_sieve_multiplicity}{online edition},
accessed 4 November 2026.

\bibitem{PNTCode}
\emph{PrimeNumberTheoremAnd}, formalization project,
\url{https://github.com/AlexKontorovich/PrimeNumberTheoremAnd},
commit \nolinkurl{c6c73610b406689c4b58325cbe1342ecca25d755};
retrieved through \url{https://github.com/lixiang90/chen_theorem},
commit \nolinkurl{11dcca8fc3ab340358503278c46d72aef072a6e7}.

\bibitem{GoldbachCode}
\emph{goldbach-lean}, large-sieve formalization,
\url{https://github.com/subfish-zhou/goldbach-lean},
commit \nolinkurl{9d4c3a4010d42aedd9e7231bf94386019cfffccc}.

\bibitem{TKS}
G. Tsoukalas, A. Kovsharov, S. Shirobokov et al.,
\emph{Advancing mathematics research with AI-driven formal
proof search},
preprint (2026),
\href{https://arxiv.org/abs/2605.22763v2}{arXiv:2605.22763v2},
Appendix~B.4.

\end{thebibliography}

\end{document}
